\documentclass[10pt,a4paper]{amsart}
\usepackage[margin=19mm]{geometry}
\usepackage{amsmath,amssymb,mathtools}
\usepackage[scr=rsfs,cal=euler]{mathalfa}
\usepackage{xfrac}
\usepackage[unicode,hidelinks,bookmarks=false]{hyperref}
\newcommand{\ie}{i.e.\ }
\newcommand{\cf}{cf.\ }
\newcommand{\resp}{respectively\ }
\newcommand{\Subsection}{\subsection}
\providecommand{\nobreakdash}{\nobreak\hbox{-}}
\setlength{\emergencystretch}{2em}
\multlinegap0pt
\usepackage{bm}
\usepackage{bbm}
\usepackage{dsfont}
\usepackage{xspace}
\usepackage{cancel}
\usepackage{mathtools}
\usepackage{amsthm}
\makeatletter
\@ifundefined{thm}{\newtheorem{thm}{Theorem}}{}
\@ifundefined{lemma}{\newtheorem{lemma}[thm]{Lemma}}{}
\makeatother
\title[Special Cases of the Shafarevich Conjecture for Complete Int]{Special Cases of the Shafarevich Conjecture for Complete Intersections in Abelian Varieties}
\author{Frank Lu}
\date{Source: arXiv:2506.14935v1 (CC BY 4.0)}
\begin{document}
\maketitle

\noindent\emph{Source and licence.} Special Cases of the Shafarevich Conjecture for Complete Intersections in Abelian Varieties. Version arXiv:2506.14935v1 (CC BY 4.0), distributed under the Creative Commons Attribution 4.0 International licence, as recorded in the arXiv licence metadata for this exact version and shown on the abstract page https://arxiv.org/abs/2506.14935v1. Full rights notice: licenses/arxiv-2506.14935-CC-BY-4.0.txt.\par\smallskip

\noindent\textit{Editorial scope.} Counted unit: the Lemma whose source label is \texttt{lem: the core q=2 inequality} in the article (source0.tex lines 510--512, source label \texttt{lem: the core q=2 inequality}), delivered together with the article's own complete original proof of it (source0.tex lines 513--527, one \texttt{proof} environment of 15 lines). No mathematics is written, repaired or re-derived: statement and proof are the article's own text line for line. Required context retained uncounted: lem: Eulerian asymptotic (lines 191--193); lem: more useable asymptotic (lines 203--209); thm: big monodromy statement (lines 372--376). Every definition, display and label of those lines is retained unchanged. Omitted from this package and not redistributed: 1--190, 194--202, 210--371, 377--509, 528--586. Nothing inside a retained excerpt is omitted or altered: every retained body is delivered exactly as the article's lines are written, so no omission receipt of any kind accompanies this package. Citations: the retained excerpts carry no citation command at all, so no bibliography entry of the article is transcribed and no reference is redistributed. Reference adaptation: none was needed and none was made; every reference inside the delivered unit resolves inside it, so no unresolved reference or citation is produced. Printed slips kept literal: no printed word, symbol, hypothesis or number of the counted statement or of its proof was altered. Layout and class: the article's own class is \texttt{article} with packages including amsmath, amsthm, amssymb, amsfonts, hyperref, geometry, graphicx, biblatex, mathrsfs, xcolor; the wrapper is the lane's pinned \texttt{amsart} editorial wrapper at 10pt on A4, which restates the theorem-like environments the retained text needs and adds the article's own macro definitions that the retained text needs (none), with extra wrapper packages (bm, bbm, dsfont, xspace, cancel, mathtools). The delivered numbering is the unit's own. Assets: no retained span contains a figure environment, a graphics-inclusion command, a PDF-page inclusion, a table or a TikZ picture; no image file is copied into this package, the package declares no asset dependency and no image file is redistributed with it. Licence: version arXiv:2506.14935v1 of this article is distributed under the Creative Commons Attribution 4.0 International licence, as recorded in the arXiv licence metadata for this exact version and shown on the abstract page https://arxiv.org/abs/2506.14935v1. A full rights notice is delivered at licenses/arxiv-2506.14935-CC-BY-4.0.txt. Characters: every retained excerpt is pure ASCII, so the delivered engine, XeLaTeX with its default Latin Modern text font, reports no missing character.
\begin{lemma}\label{lem: Eulerian asymptotic}
For all $n \geq 1, r \geq 1,$ and $k \geq 0,$ we have that $$\binom{k + r - 1}{r-1} \left((k + r)^n - (n+1)(k + r - 1)^n\right) \leq E_r(n, k) \leq \binom{k+r-1}{r-1} (k+r)^n.$$
\end{lemma}

\begin{lemma}\label{lem: more useable asymptotic}
Suppose that we are in either one of these situations:
\begin{enumerate}
    \item $n \geq 3r^2 + 100,$ $r \geq 2,$ and $k \leq 4r.$
    \item $k \leq \frac{n}{\ln(n+1) + 1} - r.$
\end{enumerate} Then, $$E_r(n, k) \geq 0.6 \binom{k+r-1}{r-1}(k+r)^n.$$ 
\end{lemma}

\begin{thm}\label{thm: big monodromy statement}
Suppose that $X$ is numerically like the complete intersection of $r$ hypersurfaces in $A,$ all representing multiples of the same ample Neron-Severi class $H,$ namely by $d_1H, d_2H, \ldots, d_rH.$ Then, if $\dim A = n \geq 10r^4 + 1000,$ there is no solution in the functions $m_H(w),$ positive integers $k$ between $2$ and $m - 1$ (where $m = \sum\limits_i m_H(i)$), and integer $s,$ to the following system of equations indexed by integers $q$ between $0$ and $n-r,$ respectively:
\[\sum\limits_{\substack{m_S: \mathbb{Z} \rightarrow \mathbb{Z} \\ 0 \leq m_S(i) \leq m_H(i) \\ \sum m_S(i) = k \\ \sum im_S(i) = s + q}} \prod\limits \binom{m_H(i)}{m_S(i)} = |\chi(X, \Omega^q_{X/K})|.\]
\par Furthermore, if the $d_i$ are all equal, there is no solution for $n \geq 10r^2 + 1000.$
\end{thm}

\begin{lemma}\label{lem: the core q=2 inequality}
For the ranges of $n$ given in Theorem \ref{thm: big monodromy statement}, we have \[\frac{1}{48} \left(\frac{\chi(X, \Omega^1_{X/K})}{\chi(X, \mathcal{O}_X)}\right)^2 > \frac{\chi(X, \Omega^2_{X/K})}{\chi(X, \mathcal{O}_X)}.\]
\end{lemma}

\begin{proof}
We start with the case that $X$ is the intersection of $r$ representatives of the same Neron-Severi class. To start, begin by employing Lemma \ref{lem: more useable asymptotic}. We then have that $\chi(X, \Omega^1_{X/K})/\chi(X, \mathcal{O}_X) > 0.6 r \left( 1 + \frac{1}{r}\right)^n$ and $\chi(X, \Omega^2_{X/K})/\chi(X, \mathcal{O}_X) < \frac{5}{3} \frac{(r+1)r}{2} \left(1 + \frac{2}{r} \right)^n.$ 
\par In other words, it thus suffices to show that \[\frac{1}{48} r^2 \left( 1 + \frac{1}{r}\right)^{2n} > \frac{125}{27} \frac{(r+1)r}{2} \left(1 + \frac{2}{r}\right)^n,\] or that \[\frac{r}{(r+1)} \left(1 + \frac{1}{r^2 + 2r}\right)^n > \frac{1000}{9}.\] Finally, using the fact that we have $r \geq 2,$ it suffices to show that $\left(1 + \frac{1}{r^2 + 2r}\right)^n > \frac{500}{3}.$
\par Now, we know that $r^2 + 2r > 6,$ and that $(1+\frac{1}{x})^x$ is increasing in positive $x;$ therefore, we can bound the left-hand side from below as $2^{\frac{n}{r^2 + 2r}}.$ But with $n > 10r^2 + 1000,$ we observe that $\frac{10r^2 + 1000}{r^2 + 2r} > 9,$ as $r^2 - 18r + 1000 > 0.$ This proves the desired inequality in the self-intersection case, since $2^9 > \frac{500}{3}.$
\par Now, suppose we are in the case where our hypersurfaces represent Neron-Severi classes $d_1H, d_2H, \ldots, d_rH.$ Without loss of generality, suppose that $d_1 \leq d_2 \leq \cdots \leq d_r.$ We start by building a lower bound for $\chi(X, \Omega^1_{X/K})$ and upper bounds for $\chi(X, \mathcal{O}_X), \chi(X, \Omega^2_{X/K}).$
\par First, for the upper bounds, we see that \begin{align*}|\chi(X, \mathcal{O}_X)| & = \frac{H^n}{n!} \sum\limits_{e \in F(r, n), |e| = n} \binom{n}{e(1), e(2), \ldots, e(r)}d_1^{e(1)} \cdots d_r^{e(r)} & = H^n \sum\limits_{e \in F(r, n), |e| = n} \prod\limits_{i=1}^r \frac{d_i^{e(i)}}{e(i)!} \\ &  = H^n \sum\limits_{e \in F(r, n), |e| = n} \prod\limits_{i=1}^r \frac{d_i}{e(i) }\frac{d_i^{e(i) - 1}}{(e(i) - 1)!} & \leq H^n \sum\limits_{e \in F(r, n), |e| = n} \prod\limits_{i=1}^r d_i \frac{d_i^{e(i) - 1}}{(e(i) - 1)!} \\ & = \frac{H^n}{(n-r)!} d_1d_2 \cdots d_r (d_1 + d_2 + \cdots + d_r)^{n-r},\end{align*} and \begin{align*}|\chi(X, \Omega^2_{X/K})| & = \frac{H^n}{n!} \sum\limits_{e \in F(r, n), |e| = n} \left(\sum\limits_{i=1}^r E(e(i), 2) + \sum\limits_{1 \leq i < j \leq r} E(e(i), 1)E(e(j), 1)\right)\binom{n}{e(1), e(2), \ldots, e(r)}d_1^{e(1)} \cdots d_r^{e(r)} \\ & = H^n \sum\limits_{e \in F(r, n), |e| = n} \left(\sum\limits_{i=1}^r E(e(i), 2) + \sum\limits_{1 \leq i < j \leq r} E(e(i), 1)E(e(j), 1)\right)\prod\limits_{i=1}^r \frac{d_i^{e(i)}}{e(i)!} \\ & \leq H^n \sum\limits_{e \in F(r, n), |e| = n} \left(\sum\limits_{i=1}^r 3^{e(i)} + \sum\limits_{1 \leq i < j \leq r} 2^{e(i) + e(j)}\right)\prod\limits_{i=1}^r \frac{d_i}{e(i) }\frac{d_i^{e(i) - 1}}{(e(i) - 1)!} \\ & \leq H^n \sum\limits_{e \in F(r, n), |e| = n} \left(\sum\limits_{i=1}^r 3^{e(i)} + \sum\limits_{1 \leq i < j \leq r} 2^{e(i) + e(j)}\right)\prod\limits_{i=1}^r d_i \frac{d_i^{e(i) - 1}}{(e(i) - 1)!} \\ & = \frac{H^n}{(n-r)!}d_1d_2 \cdots d_r \left(3\sum\limits_{i=1}^r (d_1 + d_2 + \cdots + d_r + 2d_i)^{n-r} + 4\sum\limits_{1 \leq i < j \leq r} (d_1 + d_2 + \cdots + d_r + d_i + d_j)^{n-r}\right).\end{align*}
\par Meanwhile, for the lower bound, we first claim that $$\sum\limits_{i=1}^r E(e(i), 1) \geq 0.6 \cdot 2^{e(r)}.$$ To see this, we use Lemma \ref{lem: Eulerian asymptotic}. Indeed, we know that $E(n, 1) \geq 2^n - (n+1)$ for all $n.$ If $e(r) \geq 4,$ we get the desired bound. Otherwise, $e(r) < 4;$ but then there is some $i$ for which $e(i) \geq 4,$ assuming $n > 10r^4 + 1000 \geq 4r,$ and so $0.6 \cdot 2^{e(r)} \leq 0.6 \cdot 2^{e(i)} \leq E(e(i), 1) \leq \sum\limits_{i=1}^r E(e(i), 1).$ Using this, along with the AM-GM inequality, we can then obtain the lower bound \begin{align*} |\chi(X, \Omega^1_{X/K})| & = \frac{H^n}{n!} \sum\limits_{e \in F(r, n), |e| = n} \left(\sum\limits_{i=1}^r E(e(i), 1)\right)\binom{n}{e(1), e(2), \ldots, e(r)}d_1^{e(1)} \cdots d_r^{e(r)} \\ & = H^n \sum\limits_{e \in F(r, n), |e| = n} \left(\sum\limits_{i=1}^r E(e(i), 1)\right)\prod\limits_{i=1}^r \frac{d_i^{e(i)}}{e(i)!} \\ & \geq H^n \sum\limits_{e \in F(r, n), |e| = n} \left(0.6 \cdot 2^{e(r)}\right)\prod\limits_{i=1}^r \frac{d_i}{e(i) }\frac{d_i^{e(i) - 1}}{(e(i) - 1)!} \\ & \geq H^n \sum\limits_{e \in F(r, n), |e| = n} \left(0.6 \cdot 2^{e(r)}\right)\prod\limits_{i=1}^r \frac{d_i}{(n/r)} \frac{d_i^{e(i) - 1}}{(e(i) - 1)!} \\ & = \frac{H^n}{(n-r)!}\frac{1.2 r^r}{n^r} d_1d_2 \cdots d_r \left((d_1 + d_2 + \cdots + 2d_r)^{n-r}\right).
\end{align*}
\par Using the bounds above, and cancelling out like terms, it thus suffices to show that $$\left(\frac{r^r}{n^r} (d_1 + d_2 + \cdots + 2d_r)^{(n-r)}\right)^2$$ is at least \[40 (d_1 + d_2 + \cdots + d_r)^{n-r} \left(3\sum\limits_{i=1}^r (d_1 + d_2 + \cdots + d_r + 2d_i)^{n-r} + 4\sum\limits_{1 \leq i < j \leq r} (d_1 + d_2 + \cdots + d_r + d_i + d_j)^{n-r}\right);\] this expression can be bounded above by \[40 (d_1 + d_2 + \cdots + d_r)^{n-r} \left(3r (d_1 + d_2 + \cdots + 3d_r )^{n-r} + (2r^2 - 2r) (d_1 + d_2 + \cdots + 3d_r)^{n-r}\right).\]
\par To do this, let $D = d_1 + \cdots + d_r;$ note that $D \leq rd_r.$ Then, observe that $$\frac{(d_1 + d_2 + \cdots + d_{r-1} + 2d_r)^2}{(d_1 + d_2 + \cdots + d_{r-1} + d_r)(d_1 + d_2 + \cdots + d_{r-1} + 3d_r)} = \frac{D^2 + 2Dd_r + d_r^2}{D^2 + 2Dd_r} = 1 + \frac{d_r^2}{D^2 + 2Dd_r} \geq 1 + \frac{1}{r^2 + 2r}.$$
\par Thus, to prove the lemma, it suffices to show that for the ranges of $n$ that we have given, $$\left(1 + \frac{1}{r^2 + 2r}\right)^{n-r} \geq 40\frac{2r^2 + r}{r^{2r}} n^{2r}.$$
\par Again, we observe that $\left(1 + \frac{1}{r^2 + 2r}\right)^{r^2 + 2r} \geq 2,$ so it suffices to show $2^{\frac{n-r}{r^2 + 2r}} \geq 40(2r^2 + r)\left(\frac{n}{r}\right)^{2r}$ when $n \geq 10r^4 + 1000.$ But $\frac{r}{r^2 + 2r} = \frac{1}{r+2} \leq 1/4,$ so as $2^{1/4} \leq \frac{5}{4},$ it suffices to show $2^{\frac{n}{r^2 + 2r}} \geq 50(2r^2 + r)\left(\frac{n}{r}\right)^{2r}$
\par Now, it suffices to show $2^{n/(r^2 + 2r)} \geq\left(\frac{n}{r}\right)^{2r + 2},$ since $50, 2r^2 + r \leq 10r^3 < n/r.$ Furthermore, with $r \geq 2,$ it suffices to prove the inequality with $2r + 2$ replaced with $3r.$ Taking $\log_2,$ this is equivalent to showing that $\frac{n}{r^2 + 2r} \geq 3r \log_2(\frac{n}{r})$ for $n \geq 10r^4 + 1000.$ Letting $x = \frac{n}{r},$ we want to show that $$x > 3r(r+2) \log_2x$$ for $x \geq \frac{10r^4 + 1000}{r}.$ Note that the slope of the right-hand side equals $1$ when $x = \frac{3r(r+2)}{\ln 2} \leq 10r^3,$ and furthermore that $3r(r+2) \log_2 x$ is a concave function. Therefore, if we prove the inequality holds for some $\frac{3r(r+2)}{\ln 2} \leq x < \frac{10r^4 + 1000}{r},$ then it will hold for all larger $x$ (since $3r(r+2) \log_2 x$ has slope less than $1$ for $x \geq \frac{3r(r+2)}{\ln 2}$). 
\par But now at $x = 10r^3,$ we have that $\frac{10r^2}{3(r+2)} > \log_2 10r^3$ for $r \geq 3,$ by using a similar argument to the above and a bit of direct computation. For $r = 2,$ we can check that $x = \frac{10r^4 + 1000}{r}$ satisfies the desired inequality, and so we are done.
\end{proof}

\end{document}
